% Arquivo: latex/metodo.tex
% Seção 3: Método
% Gerado a partir do roteiro latex/roteiro_metodo.md

\subsection{Estratégia de busca e seleção}
\label{subsec:estrategia-busca}

O objetivo desta revisão sistemática é identificar e sintetizar estudos empíricos que
empreguem abordagem econométrica para avaliar o efeito de instrumentos de financiamento
da PNDR --- Fundos Constitucionais de Financiamento, Fundos de Desenvolvimento Regional
e Incentivos Fiscais --- sobre indicadores socioeconômicos locais. O protocolo de busca,
triagem e seleção segue as diretrizes PRISMA 2020 (\textit{Preferred Reporting Items for
Systematic Reviews and Meta-Analyses}), conforme \citeonline{Pageetal2021}, com o intuito
de assegurar transparência e reprodutibilidade em cada etapa do processo.

A busca foi conduzida em cinco bases de dados acadêmicas --- Scopus, SciELO, Portal de
Periódicos da CAPES, EconPapers/RePEc e anais dos encontros da ANPEC. A consulta a
múltiplas fontes justifica-se pela elevada dispersão dos registros: periódicos nacionais
de economia tendem a apresentar menor cobertura em bases de dados internacionais, e
parte relevante da produção sobre instrumentos da PNDR concentra-se em \textit{working
papers}, textos para discussão e apresentações em congressos nacionais, veículos nem
sempre indexados por essas bases. Das cinco fontes, obtiveram-se 137 registros. Após a
remoção de 19 duplicatas em três fases automatizadas, 118 registros únicos foram
submetidos à análise de elegibilidade. A triagem combinou classificação automatizada com
revisão dos autores, resultando na inclusão de 45 estudos e exclusão de
73. A Figura~\ref{fig:diagrama_prisma} apresenta o diagrama de fluxo PRISMA 2020 com
o detalhamento quantitativo de cada etapa.\footnote{Os scripts de coleta, deduplicação
e triagem estão disponíveis no repositório de replicação do autor em
\url{https://github.com/hermelino/pndr_survey}.}

% TODO: Inserir diagrama PRISMA 2020 com números atualizados
% Fluxo: 5 bases → 137 registros → dedup (19 removidas) → 118 únicos
%        → triagem → 45 aprovados, 73 rejeitados

\subsection{Bases de dados consultadas}
\label{subsec:bases-dados}

A seleção das cinco bases visa garantir cobertura abrangente da literatura. A
Scopus indexa os principais periódicos internacionais de economia; a SciELO
complementa a cobertura de periódicos nacionais em acesso aberto; o Portal de
Periódicos da CAPES agrega centenas de bases de dados nacionais e internacionais;
o EconPapers/RePEc concentra \textit{working papers} e textos para discussão com
ampla cobertura em economia; e os anais da ANPEC capturam trabalhos apresentados
nos encontros nacionais e regionais de economia, incluindo estudos ainda não
publicados em periódicos. A Tabela~\ref{tab:bases-registros} resume o número de
registros obtidos em cada fonte.

\begin{table}[h]
\centering
\caption{Registros obtidos por base de dados}
\label{tab:bases-registros}
\begin{tabular}{lr}
\toprule
Base de dados & Registros \\
\midrule
Scopus & 16 \\
SciELO & 5 \\
Portal de Periódicos da CAPES & 30 \\
EconPapers/RePEc & 24 \\
ANPEC (encontros nacionais e regionais) & 62 \\
\midrule
\textbf{Total bruto} & \textbf{137} \\
\bottomrule
\end{tabular}
\fonte{Elaboração própria.}
\end{table}

\subsection{Expressões de busca e critérios de inclusão e exclusão}
\label{subsec:criterios}

A estratégia de busca empregou combinação de três blocos de termos conectados pelo
operador AND: (1)~instrumentos de financiamento (``fundo constitucional'', ``fundo
de desenvolvimento'', ``incentivo fiscal'', com variações em inglês); (2)~organizações
e siglas associadas à PNDR (FNE, FNO, FCO, FDNE, FDCO, SUDENE, SUDAM, PNDR); e
(3)~delimitação geográfica (Nordeste, Amazônia, Centro-Oeste, Brasil). Dentro de
cada bloco, os termos foram conectados pelo operador OR. As expressões exatas foram
adaptadas à sintaxe de cada base.

São admissíveis a compor a amostra artigos publicados em periódicos, \textit{working
papers}, textos para discussão e apresentações em congressos. Os critérios de inclusão
adotados foram: (i)~estudos empíricos que avaliem pelo menos um instrumento de
financiamento da PNDR; (ii)~emprego de abordagem econométrica que permita isolamento
do efeito da política sobre indicadores socioeconômicos; e (iii)~publicação em
qualquer dos formatos acima.

Os critérios de exclusão foram: (a)~estudos que não avaliem explicitamente
instrumentos da PNDR; (b)~abordagens exclusivamente qualitativas ou descritivas;
(c)~avaliações sem emprego de métodos econométricos de controle de fatores simultâneos
à política; (d)~estudos anteriores a 2005, ano a partir do qual se identificam as
primeiras avaliações empíricas desses instrumentos; e (e)~duplicatas de versão
publicada, casos em que um mesmo trabalho aparece como texto para discussão ou
apresentação em congresso e como artigo em periódico, retendo-se a versão publicada.

\subsection{Processo de deduplicação}
\label{subsec:deduplicacao}

A consulta a múltiplas bases de dados implica a presença de registros duplicados,
uma vez que um mesmo estudo pode estar indexado em mais de uma fonte. Para tratamento
dessas ocorrências, aplicou-se algoritmo automatizado em três fases sequenciais:
correspondência por identificador DOI, similaridade de título e identidade de arquivo
PDF. Ao todo, 19 duplicatas foram identificadas e removidas, reduzindo o conjunto de
137 para 118 registros únicos. Em caso de duplicata entre bases, adotou-se como
critério de retenção a seguinte ordem de prioridade: Scopus, SciELO, Portal CAPES,
EconPapers e ANPEC, privilegiando a versão com metadados mais completos.

\subsection{Triagem final}
\label{subsec:triagem-final}

Os 118 registros únicos foram submetidos a processo de triagem em duas etapas. Na
primeira, uma classificação automatizada identificou os registros com maior
probabilidade de atender aos critérios de inclusão. Na segunda, os registros foram
avaliados integralmente, com revisão das classificações automatizadas. A decisão
final de inclusão ou exclusão baseou-se na avaliação dos autores.

Dos 118 registros, 45 foram incluídos na revisão e 73 excluídos. A
Tabela~\ref{tab:triagem} apresenta a distribuição dos motivos de exclusão.

\begin{table}[h]
\centering
\caption{Registros excluídos por motivo de rejeição}
\label{tab:triagem}
\begin{tabular}{lr}
\toprule
Motivo de exclusão & Registros \\
\midrule
Sem instrumentos da PNDR & 35 \\
Sem método econométrico & 17 \\
Documento não científico & 10 \\
Duplicata de versão publicada & 7 \\
Anterior a 2005 & 2 \\
Outros & 2 \\
\midrule
\textbf{Total excluídos} & \textbf{73} \\
\bottomrule
\end{tabular}
\fonte{Elaboração própria.}
\end{table}

\subsection{Descrição dos estudos obtidos}
\label{subsec:descricao-estudos}

A aplicação dos critérios descritos nas subseções anteriores resultou em amostra
de 45 estudos empíricos, datados de 2005 a 2025. Destes, 22 são artigos publicados
em periódicos nacionais ou internacionais, e os demais 23 são textos para discussão
ou trabalhos apresentados em congressos.

A produção acadêmica sobre o tema distribui-se de forma desigual ao longo do
período. De 2005 a 2013, identificaram-se apenas 8 estudos. No intervalo de 2014 a
2019, a produção elevou-se para 11 trabalhos, refletindo o amadurecimento da
literatura. O crescimento mais expressivo concentrou-se entre 2020 e 2025, com 26
estudos --- mais da metade da amostra ---, indicando intensificação recente do
interesse acadêmico pela avaliação empírica dos instrumentos da PNDR.

A escala \textit{Maryland Scientific Methods} (MSM), descrita em
\citeonline{MadalenoWaights2016}, classifica a robustez dos métodos de avaliação de
impacto em escores de 1 (menos robusto) a 5 (mais robusto), segundo a capacidade de
lidar com vieses de seleção. Dentre os estudos selecionados, predominam abordagens
com escore 3 --- métodos que controlam por características observáveis, porém sem
garantir completa eliminação de vieses de seleção. Os métodos mais frequentes
incluem painel de efeitos fixos, diferenças em diferenças, \textit{Propensity Score
Matching}, painel espacial e regressão MQO. Abordagens quase experimentais com
escore 4, como \textit{Regression Discontinuity Design}, permanecem pouco
utilizadas. A ausência de avaliações com escore máximo de robustez reflete as
restrições inerentes ao desenho não aleatório da política, em que a alocação de
recursos tende a ser determinada pela demanda das firmas.

% TODO: Inserir figura do mapa de distribuição dos artigos (litmap)
% atualizado com 45 estudos
